% Fuente conservada íntegra: originales/cuaderno/NOTA.md, §§1--7.
\section{Mémoire chronologique, bilan et réponse des vacances}
\label{md:cron:seccion}

La construction suivante reçoit le facteur symbolique de capacité et étudie ses lecteurs ultérieurs. Le rapport interne $3^6/10^3=729/1000$ et son calendrier appartiennent à la construction précédente (sections~\ref{mono-ampl:relojes} et~\ref{vac-capacidad-trit}). L'architecture APP--TRIT--TPK, les régions, le retour nonadique et le registre $K$ conservent leur position causale antérieure. Les logarithmes employés dans le lecteur de capacité décrivent ce facteur et ne sélectionnent pas à nouveau les régions de $\pi,\varphi,e,\alpha$. Les propositions sont des conséquences focales concernant l'observation et la mémoire ; leur portée se distingue de l'état TPK complet et de l'attribution d'une priorité historique à des résultats généraux sur les suites mécaniques ou l'information.

\subsection{Objet, mesure et deux lecteurs}
\label{md:cron:lectores}

Soit $\delta=\log_{10}(10/9)$ la pente du calendrier de capacité, et soit $\theta$ une phase uniformément distribuée dans $[0,1)$. La mesure représente l'incertitude sur l'origine de lecture, tandis que le mécanisme demeure déterminé. Nous définissons
\begin{equation}
 \begin{aligned}
 z_j(\theta)&=\lfloor(j+1)\delta+\theta\rfloor
             -\lfloor j\delta+\theta\rfloor,\\
 W_n(\theta)&=(z_0(\theta),\ldots,z_{n-1}(\theta)),\\
 B_n(\theta)&=\sum_{j=0}^{n-1}z_j(\theta).
 \end{aligned}
 \label{md:cron:dos-lectores}
\end{equation}
Le bloc $W_n$ conserve la chronologie et $B_n$ conserve le bilan. Pour une phase exacte et une règle exacte connues, les deux objets sont déterministes ; les entropies se rapportent expressément à l'ensemble des phases possibles.

La sommation télescopique donne
\begin{equation}
 B_n(\theta)=\lfloor n\delta+\theta\rfloor.
 \label{md:cron:telescopia}
\end{equation}
Ainsi, $B_n$ n'admet que les valeurs $\lfloor n\delta\rfloor$ et $\lceil n\delta\rceil$. Les points $\{-j\delta\}$, avec $j=0,\ldots,n$, divisent le cercle en $n+1$ intervalles de phase. Chaque intervalle porte une suite distincte de longueur $n$. En effet, les préfixes cumulés $\lfloor k\delta+\theta\rfloor$, pour $k=1,\ldots,n$, sont monotones en $\theta$, et chaque frontière modifie l'un d'eux ; la suite reconstruit tous ces préfixes.

Les extrémités sont de mesure nulle. Une convention semi-ouverte résout leur assignation et conserve les entropies.

\subsection{Information croissante et taux entropique nul}
\label{md:cron:crecimiento}

Soient $\ell_0,\ldots,\ell_n$ les longueurs des intervalles de phase et
\begin{equation}
 H_n=-\sum_{j=0}^{n}\ell_j\log_2\ell_j.
 \label{md:cron:entropia-bloque}
\end{equation}

\begin{proposition}
\label{md:cron:proposicion-tasa}
On a simultanément
\begin{equation}
 H_n\longrightarrow+\infty,\qquad
 \frac{H_n}{n}\longrightarrow0.
 \label{md:cron:dos-limites}
\end{equation}
\end{proposition}
\begin{proof}
La borne supérieure est $H_n\leq\log_2(n+1)$. Si $\ell_{\max}(n)=\max_j\ell_j$, alors $H_n\geq-\log_2\ell_{\max}(n)$. La pente $\delta$ est irrationnelle : si $\delta=p/q$ avec $q>0$, on aurait $(10/9)^q=10^p$, ce qui contredit la valuation du nombre premier $3$. Par irrationalité, les points de l'orbite sont denses. Les partitions successives ont un pas maximal tendant vers zéro : pour chaque $\epsilon>0$, il existe un segment fini $\epsilon$-dense, et les segments ultérieurs ne font que subdiviser ses intervalles. La borne inférieure diverge donc. La borne supérieure divisée par $n$ tend vers zéro.
\end{proof}

Les longueurs des intervalles ne sont pas égales en général. La preuve établit les bornes précédentes sans affirmer $H_n=\log_2(n+1)$ ni une loi asymptotique logarithmique exacte qui requerrait des conditions diophantiennes supplémentaires. L'entropie topologique du facteur est également nulle parce que $\log(n+1)/n\to0$. L'information du bloc et son taux sont des quantités différentes. La croissance de l'archive d'une trajectoire connue n'équivaut pas, à elle seule, à une production d'information statistique indépendante.

\subsection{Information chronologique non déterminée par le bilan}
\label{md:cron:condicional}

\begin{corollary}
\label{md:cron:corolario-condicional}
Pour le même ensemble de phases,
\begin{equation}
 H(W_n\mid B_n)=H_n-H(B_n)
 \geq H_n-1\longrightarrow+\infty.
 \label{md:cron:entropia-condicional}
\end{equation}
\end{corollary}
\begin{proof}
$B_n$ est une fonction de $W_n$, de sorte que la règle de chaîne donne l'égalité. Son support possède deux éléments, d'où $H(B_n)\leq1$ bit. On applique la proposition~\ref{md:cron:proposicion-tasa}.
\end{proof}

L'entropie conditionnelle divergente est la moyenne sur les bilans ; on n'attribue pas la même incertitude à chaque bilan concret. Cette perte n'est pas non plus attribuée à l'état enrichi complet : lorsque la mémoire retenue restitue $W_n$, on a
\begin{equation}
 H(W_n\mid B_n,M_n)=0.
 \label{md:cron:restitucion-memoria}
\end{equation}
Par cardinalité, au moins une des deux fibres de $B_n$ contient $\lceil(n+1)/2\rceil$ suites ou davantage. Cette borne combinatoire se distingue de la borne entropique. Le résultat précise la relation entre valeur et généalogie dans ce lecteur : les deux bilans possibles n'épuisent pas les historiques temporels compatibles.

\subsection{Courant centré : bilan borné et activité persistante}
\label{md:cron:corriente}

Le lecteur centré s'exprime par
\begin{equation}
 J_j=z_j-\delta,\qquad
 \sum_{j=0}^{n-1}J_j=B_n-n\delta.
 \label{md:cron:corriente-definicion}
\end{equation}
Par conséquent, $\left|\sum_{j=0}^{n-1}J_j\right|<1$ pour toute fenêtre. Simultanément,
\begin{equation}
 \lim_{n\to\infty}\frac1n\sum_{j=0}^{n-1}J_j=0,\qquad
 \lim_{n\to\infty}\frac1n\sum_{j=0}^{n-1}J_j^2
 =\delta(1-\delta)>0.
 \label{md:cron:actividad}
\end{equation}
\begin{proof}
On a $z_j^2=z_j$ et $B_n/n\to\delta$ par la borne de discrépance. Le développement $(z_j-\delta)^2=z_j(1-2\delta)+\delta^2$, suivi de la substitution de la moyenne de $z_j$, donne la seconde identité. L'argument vaut pour toute phase et ne requiert pas de moyenne aléatoire.
\end{proof}

Le bilan accumulé borné, l'activité quadratique positive et le taux d'entropie nul sont donc compatibles. Les identités conservent la distinction entre le lecteur centré et une représentation de charge élémentaire ; elles n'établissent pas non plus de production énergétique sans bilan.

Si une réalisation dimensionnelle déjà établie prend $I_j=(u_Q/t_0)J_j$, on obtient
\begin{equation}
 I_{\mathrm{rms}}^2
 =\left(\frac{u_Q}{t_0}\right)^2\delta(1-\delta).
 \label{md:cron:eficaz}
\end{equation}
Si chaque $I_j$ est maintenu pendant un intervalle de durée $t_0$, cette moyenne discrète coïncide avec la moyenne temporelle. À titre de comparaison physique ultérieure, une charge résistive idéale $R_{\mathrm{res}}>0$ aurait pour puissance moyenne
\begin{equation}
 \overline P_{\mathrm{res}}
 =R_{\mathrm{res}}\left(\frac{u_Q}{t_0}\right)^2
 \delta(1-\delta)>0.
 \label{md:cron:potencia-resistiva}
\end{equation}
C'est une prédiction du modèle composé avec cette charge. La résistance et le dispositif conservent leur propre spécification ; ils ne se dérivent pas de la seule suite de capacité.

\subsection{Spectre du courant et orientation}
\label{md:cron:espectro}

La fonction d'observation est $f(\theta)=\mathbf1_{[1-\delta,1)}(\theta)-\delta$, avec $J_j=f(\theta+j\delta)$. Ses coefficients de Fourier, pour $m\neq0$, sont
\begin{equation}
 a_m=\frac{\exp(2\pi i m\delta)-1}{2\pi i m},\qquad
 |a_m|^2=\frac{\sin^2(\pi m\delta)}{\pi^2m^2},\qquad
 a_0=0.
 \label{md:cron:fourier}
\end{equation}
Ils s'obtiennent en intégrant l'indicatrice sur son intervalle. L'identité de Parseval donne
\begin{equation}
 \sum_{m\neq0}|a_m|^2=\delta(1-\delta),
 \label{md:cron:parseval}
\end{equation}
en accord avec l'activité quadratique. Les fréquences temporelles sont $m\delta$ modulo $1$. La corrélation stationnaire, avec phase uniforme, est déterminée par ces intensités ; celles-ci ne contiennent pas la phase initiale. Une translation multiplie $a_m$ par une phase unitaire, et l'inversion temporelle échange les fréquences opposées.

La puissance spectrale isolée ne restitue ni l'orientation ni l'origine temporelle. Une lecture sensible à la phase ou une référence marquée apporte l'information complémentaire. La déduction conserve la distinction entre cette référence générale et le registre $K$, ainsi qu'entre les rythmes symboliques et une réalisation expérimentale déterminée.

\subsubsection{Référence cyclique et transports de même intensité spectrale}
\label{md:cron:k-ciclico}

La propriété cyclique du registre $K$ permet une conséquence finie
indépendante du spectre précédent. Dans $V_K=\mathbb R^{12}$, soit $S$ le
décalage $(Sx)_i=x_{i+1}$ et
\begin{equation}
 O_K=(K,SK,\ldots,S^{11}K).
 \label{md:cron:matriz-ciclica}
\end{equation}
On utilise l’inversibilité de $O_K$ établie par la référence cyclique
du registre. Pour tout opérateur réel $H$ qui commute avec $S$,
posons $y_+=HK$ et $y_-=H^TK$.

\begin{proposition}
\label{md:cron:proposicion-intensidades}
Les douze intensités de Fourier de $y_+$ et $y_-$ coïncident.
Cependant, $y_+=y_-$ si et seulement si $H=H^T$.
Chaque réponse vectorielle complète, avec $K$ et l’orientation de $S$,
retrouve son opérateur.
\end{proposition}
\begin{proof}
$H$ est circulant. Ses multiplicateurs de Fourier sont $\lambda_m$ et ceux
de $H^T$ sont leurs conjugués, de sorte que les deux réponses ont
les intensités $|\lambda_m|^2|\widehat K_m|^2$.
Si leurs vecteurs coïncident, $(H-H^T)K=0$. La commutation implique que
$H-H^T$ annule également $S^jK$ pour tout $j$. Ces vecteurs forment une base ;
par conséquent, $H-H^T=0$. L’implication inverse est immédiate.
La récupération est exprimée par
\begin{equation}
 H=O_{y_+}O_K^{-1},\qquad
 H^T=O_{y_-}O_K^{-1}.
 \label{md:cron:restitucion-operador}
\end{equation}
\end{proof}

En particulier, les décalages $S$ et $S^{-1}$ produisent des réponses de
même spectre de puissance, mais distinctes. Les autocorrélations exactes
du registre donnent
\begin{equation}
 \|SK-S^{-1}K\|^2
 =2\bigl(\|K\|^2-\langle K,S^2K\rangle\bigr)
 =2175744>0.
 \label{md:cron:separacion-desplazamientos}
\end{equation}
L’inversion rationnelle de $O_K$ retrouve à partir des deux réponses les
coefficients $e_1$ et $e_{11}$, respectivement. $K$ est utilisé après
sa génération comme référence marquée. L’intensité spectrale ou le
bilan scalaire ne remplacent pas la réponse vectorielle.
Le corollaire décrit le support cyclique à douze composantes ;
il conserve la différence entre $S$ et la dynamique enrichie complète,
ainsi que les opérateurs supplémentaires d’une conjugaison physique de charge.
Pour un $H$ général, la transposition ne signifie pas non plus une inversion temporelle ;
dans le cas $S$, elle coïncide bien avec $S^{-1}$.

\subsection{Lien opératoriel électronique}
\label{md:cron:electron}

L’incidence entre les fibres de somme et de produit d’APP sur $D_9^3$
définit $C$. Son vecteur tritique
$m_{\mathrm{ph}}=(1,-1,0)^3$ satisfait
\begin{equation}
 Cm_{\mathrm{ph}}=C^Tm_{\mathrm{ph}}
 =-\frac12m_{\mathrm{ph}}.
 \label{md:cron:modo-tritico}
\end{equation}
Le calendrier induit parcourt les régimes du même sélecteur tritique ;
l’incidence sélectionne le plan principal de singularité $1/2$.
Dans celui-ci,
\begin{equation}
 P=\begin{pmatrix}1&0\\0&0\end{pmatrix},\qquad
 Q=\begin{pmatrix}1/4&\sqrt3/4\\\sqrt3/4&3/4\end{pmatrix},
 \qquad J=\frac4{\sqrt3}[P,Q],\qquad J^2=-I.
 \label{md:cron:plano-electronico}
\end{equation}
L’incidence et la norme du commutateur sont développées dans la
section~\ref{vi:dep:prefactor}. On retrouve ainsi une relation effective entre
l’organisation temporelle et le bloc électronique, en conservant le sélecteur,
ses étiquettes et l’incidence. La représentation de charge et la réalisation
dimensionnelle maintiennent leurs opérateurs propres : le lien n’assigne pas
directement les charges élémentaires opposées à $z_j-\delta$.

\subsection{Information retenue et composition énergétique}
\label{md:cron:informacion-energia}

La récupération de fibre établit
\begin{equation}
 H(X\mid q(X),M(X))=0
 \label{md:cron:inversa-enriquecida}
\end{equation}
lorsque le lecteur enrichi a un inverse ; une mise à jour bijective
conserve $H(X)$. Ces propriétés sont réunies dans
\eqref{eq:ii-kb-aut-fibra} et~\eqref{eq:ii-kb-aut-reversible}.
La coexistence de $H_n$ croissant et d’un taux nul est compatible avec ces
propriétés et avec les identités de Shannon.

Le travail d’effacement dépend de l’information accessible retenue.
Dans le régime idéal isotherme des mémoires logiques dégénérées et dans la
limite asymptotique appropriée, le coût par copie avec information latérale est
\begin{equation}
 W_{\mathrm{borrado,ideal}}
 =k_BT\ln(2)\,H(X\mid M).
 \label{md:cron:borrado}
\end{equation}
Si $X$ est retrouvé à partir de $M$, la contribution logique idéale peut être nulle.
Le transport physique, la lecture, le contrôle et le cycle complet conservent
des bilans propres ; l’égalité $h_{\mathrm{top}}=0$ ne les élimine pas.

La composition pertinente conserve la séquence
\[
 \begin{gathered}
 \text{chronologie}\longrightarrow
 \text{registre de frontière}\longrightarrow
 \text{lecteur de courant et d’état}\\
 \longrightarrow
 \text{mémoire conservée}\longrightarrow
 \text{bilan thermique}.
 \end{gathered}
\]
Le développement précédent construit les deux premiers lecteurs et leurs
conséquences informationnelles. Sa portée est distinguée de la certification
d’un dispositif physique complet et maintient intacte la généalogie
constructive des constantes.
